\definesection{Line Directives}
In order for a \g{directive} to be a \g{line directive}, its \g{identifier} must \g{match} the beginning of a \g{line}. \\

A \g{textual component} specified by a \g{line directive} can potentially contain any other \g{textual component}. Therefore, any \g{directive} is potentially recognisable within a \g{line directive}, including other \gpl{line directive}. However, a \g{line directive} may explicitly restrict which \gpl{directive} are recognised within itself. A \g{line directive} cannot cross the boundaries of another \g{line directive} of a different kind. If such a case were to occur, the current \g{line directive} is forcibly ended without regard for any possible trailing \g{match}.

\definesubsubsection{Singular Line Directives}
A \g{line directive} is a \g{singular line directive} if it is only ever active for a single \g{line}. If it is matched on two consecutive \gpl{line} this results in two separate \gpl{resulting textual component}. \\

When a \g{singular line directive} is \glink{processing}{processed}, processing begins anew over the \g{content binding} until the end of the \g{line} is reached, at which point the \g{resulting textual component} is ended. After that, control is handed back to the \g{standard processing loop}.

\definesubsubsection{Spanning Line Directives}
A \g{line directive} is a \g{spanning line directive} if the \g{identifier} contains a \g{content binding}, and if \glink{match}{matches} on consecutive \gpl{line} of the \g{identifier} are interpreted as a single \g{match}. The semantics of such a spanning match are as follows: Only a single \g{resulting textual component} is produced for all the consecutively \glink{match}{matching} \gpl{line}. The \g{text} of this \g{resulting textual component} is produced by concatenating the contents of the \g{content binding} on each \g{line}. If the \g{content binding} does not \g{match} the \g{newline} on every \g{line}, the \g{newline} must be inserted between each \g{string} of the \g{content binding}. \\

When a \g{spanning line directive} is \glink{processing}{processed}, processing begins anew over the \g{content binding} until the end of the \g{line} is reached. Standard end of \g{line} \g{interpretation} proceeds. If the following \g{line} \glink{match}{matches} the same \g{spanning line directive} as before, processing begins anew over the \g{content binding} thereof without any new \gpl{resulting textual component} being started or inserted. If the following \g{line} does not \g{match} the same \g{spanning line directive} as before, the \g{resulting textual component} is ended and control is handed back to the \g{standard processing loop}.

\definesubsubsection{Guarded Line Directives}
A \g{line directive} is a \g{guarded line directive} if its \glink{match}{matched} region is specified by two \gpl{identifier} that each match a single \g{line}. The \g{text} of the \g{resulting textual component} is the \g{text} from the \g{line} immediately after the \g{line} the first \g{identifier} \glink{match}{matches} until and including the \g{line} immediately before the \g{line} the second \g{identifier} \glink{match}{matches}. \\

When a \g{guarded line directive} is \glink{processing}{processed}, processing begins anew over the \g{content binding} until the the part of the \g{identifier} after the \g{content binding} is \glink{full match}{fully matched}, at which point the \g{resulting textual component} is ended and control is handed back to the \g{standard processing loop}.

\input{5.1-paragraph.tex}
\definesubsection{Blockquote}
\begin{identifier}{blockquote header}
\~ <content !(\| )+>(\| <body .*>)?
\end{identifier}
\begin{identifier}{blockquote body}
\| <content .*>
\end{identifier}
\definetextualcomponent{blockquote header}{margin: left; font-weight: bold}
\definetextualcomponent{blockquote body}{margin: left} \\

The blockquote header is a \g{singular line directive} that identifies the source of a quote. Only the \g{text} held by the \g{content binding} is outputted into the \g{resulting textual component}. The blockquote header \g{content binding} can only contain \gpl{inline directive}. If the \inline$body$ \g{binding} is present, then the blockquote header's \g{resulting textual component} is closed, and parsing resumes as if at a block toplevel, meaning a blockquote body will be parsed next. This allows a shorter form of combined header and body. \\

The blockquote body is a \g{spanning line directive} that identifies a body of \g{text} that is being quoted. The blockquote body can contain any \g{directive} with the condition that the \gpl{directive} are matched against the \g{text} of the \g{resulting textual component}. \\

An implementation may choose to group the \comp{blockquote header} and \comp{blockquote body} together and reorder them if they are found consecutive to one another. However, a body can only ever be grouped together with a single header. In the case where a header lies between two bodies, the header is counted to belong to the second body. If a header is found without a corresponding body, the \g{implementation} may \glink{signalling}{signal} a \g{warning}. \\

\begin{examples}
  \begin{examplesource}
~ This Document
| The blockquote header is a \
| singular line directive.
  \end{examplesource}
  \begin{exampleoutput}
    \begin{blockquote}[This Document]
      The blockquote header is a singular line directive.
    \end{blockquote}
  \end{exampleoutput}
  \begin{examplesource}
| Unattributed text.
  \end{examplesource}
  \begin{exampleoutput}
    \begin{blockquote}
      Unattributed text.
    \end{blockquote}
  \end{exampleoutput}
  \begin{examplesource}
~ Yukari | Hello there!
  \end{examplesource}
  \begin{exampleoutput}
    \begin{blockquote}[Yukari]
      Hello there!
    \end{blockquote}
  \end{exampleoutput}
\end{examples}

%%% Local Variables:
%%% mode: latex
%%% TeX-master: "markless"
%%% TeX-engine: luatex
%%% TeX-command-extra-options: "-shell-escape"
%%% End:

\definesubsection{Lists}
\begin{identifier}{ordered-list}
<number ~d+>\.<content .*>
(<spacing ~_+> <content .*>)*
\end{identifier}
\begin{identifier}{unordered-list}
- <content .*>
(<spacing ~_+> <content .*>)*
\end{identifier}
\definetextualcomponent{ordered list}{margin: left}
\definetextualcomponent{ordered list item}{display: list-item; list-item-prefix: number}
\definetextualcomponent{unordered list}{margin: left}
\definetextualcomponent{unordered list item}{display: list-item; list-item-prefix: dot} \\

The lists are \gpl{spanning line directive} and mark the enumeration of one or more items of a list. They can contain any \g{directive} with the condition that the \gpl{directive} are matched against the \g{text} of the \g{resulting textual component}. \\

After the respective list \g{identifier} has been \glink{match}{matched}, a new respective item \g{textual component} in which the higher \g{level} \g{text} is contained, is inserted for each \g{match} into the spanning \g{resulting textual component}. A single \g{match} may span over multiple \gpl{line} if the \g{text} \glink{match}{matched} by the \inline$spacing$ \g{binding} is of the same length as that of the \inline$number$ \g{binding}. In such a case, each item \g{match} itself is treated like a \g{spanning line directive} where the \g{content binding} is concatenated. \\

Ordered list items must be numbered by the \g{decimal number} given by the \inline$number$ \g{binding}, even if there is no order to how the numbers appear in the list or if there are duplicates. \\

\begin{examples}
\begin{examplesource}
- Finish this spec
- Implement a parser
\end{examplesource}
  \begin{exampleoutput}
    \begin{minipage}{0.5\textwidth}
      \begin{itemize}[noitemsep]
      \item Finish this spec
      \item Implement a parser
      \end{itemize}
    \end{minipage}
  \end{exampleoutput}
\begin{examplesource}
1.Buy some ingredients
2.Clean the kitchen
  Don't forget the sink!
5.Watch TV
\end{examplesource}
  \begin{exampleoutput}
    \begin{minipage}{0.5\textwidth}
      \begin{enumerate}[noitemsep]
      \item Buy some ingredients
      \item Clean the kitchen\\Don't forget the sink!
        \setcounter{enumi}{4}
      \item Watch TV
      \end{enumerate}
    \end{minipage}
  \end{exampleoutput}
\end{examples}

%%% Local Variables:
%%% mode: latex
%%% TeX-master: "markless"
%%% TeX-engine: luatex
%%% TeX-command-extra-options: "-shell-escape"
%%% End:

\input{5.4-header.tex}
\input{5.5-horizontal-rule.tex}
\definesubsection{Code block}
\begin{identifier}{code-block}
<prefix ::+> *<language ![,]+>?<options .*>
<content .*>
<prefix>
\end{identifier}
\definetextualcomponent{code block}{font-family: monospace; white-space: preserve} \\

The code block is a \g{guarded line directive}. It marks the \g{text} to belong to a \g{textual component} that somehow distinguishes the block as source code. Only the \g{text} held by the \g{content binding} is outputted to the \g{resulting textual component}. The code block \g{directive} cannot contain any other \gpl{directive}. \\

The \gpl{newline} and \g{whitespace} must be represented exactly as in the source text. Multiple consecutive \g{whitespace} \gpl{character} cannot be combined and must be individually represented. A \g{newline} \g{character} cannot be escaped and must always result in a new line being started. \G{escaping} is deactivated within the content, meaning backslashes are output literally in the \g{resulting textual component}.  \\

The \inline$options$ \g{binding} holds potential parameters that can configure the \g{style} of the \g{resulting textual component}. The syntax and effect of the options is \g{implementation dependant}. If a language is requested in the \linline$language$ \g{binding} that the \g{implementation} does not have specific support for, a \g{warning} is \glink{signalling}{signalled}. The language \inline$text$ must always be supported and will incur no styling changes. \\

\begin{examples}
\begin{examplesource}
Some unexciting code:
:: common-lisp
(print "Hello world")
::
\end{examplesource}
  &$\Rightarrow$&
  \begin{tabular}{@{}l@{}}
Some unexciting code: \\
\begin{lstlisting}[style=codestyle,language=Lisp,showstringspaces=false]
(print "Hello world")
\end{lstlisting}
\end{tabular}
\end{examples}

%%% Local Variables:
%%% mode: latex
%%% TeX-master: "markless"
%%% TeX-engine: luatex
%%% TeX-command-extra-options: "-shell-escape"
%%% End:

\definesubsection{Instruction}
\begin{identifier}{instruction}
\! <instruction .*>
\end{identifier}

The instruction is a \g{singular line directive}. Its purpose is to interact with the \g{implementation} and cause it to perform differently. There is no corresponding \g{resulting textual component} for the instruction \g{directive}. \\

An \g{implementation} is allowed to add further instructions. If an instruction is not recognised, the \g{implementation} must \glink{signalling}{signal} an \g{error}. \\

\input{5.7.1-set.tex}
\input{5.7.2-message.tex}
\input{5.7.3-include.tex}
\definesubsubsection{Directives}
\defineinstruction{disable}{disable <directive ![ ]+>( <directive ![ ]+>)*}
\defineinstruction{enable}{enable <directive ![ ]+>( <directive ![ ]+>)*} \\

The \instr{disable} and \instr{enable} instructions cause the \g{implementation} to respectively disable or enable the named \gpl{directive}. If a given name is not recognised, the \g{implementation} may \glink{signalling}{signal} a \g{warning}. If a user tries to disable the \comp{paragraph} \g{directive}, an \g{error} must be \glink{signalling}{signalled}. \\

\begin{examples}
\begin{examplesource}
! disable instruction
! error Exit!
\end{examplesource}
  \begin{exampleoutput}
    ! error Exit!
  \end{exampleoutput}
\end{examples}

%%% Local Variables:
%%% mode: latex
%%% TeX-master: "markless"
%%% TeX-engine: luatex
%%% TeX-command-extra-options: "-shell-escape"
%%% End:

\definesubsubsection{Label}
\defineinstruction{label}{label <name .*>} \\

Associates a \g{label} with the \g{component} that immediately precedes this \g{instruction} at the current level. The \g{implementation} must \glink{signalling}{signal} an \g{error} if there is no preceding component at the current level. \\

%%% Local Variables:
%%% mode: latex
%%% TeX-master: "markless"
%%% TeX-engine: luatex
%%% TeX-command-extra-options: "-shell-escape"
%%% End:

\definesubsubsection{Raw}
\defineinstruction{raw}{raw <backend ![ ]+> <content .*>} \\

If the \g{implementation}'s chosen output backend matches that of the \inline$backend$ \g{binding}, the \g{implementation} should emit the \inline$content$ \g{binding}'s text verbatim into the resulting document. This should allow creating output specific effects. The exact semantics and results of this are \g{implementation dependant}. If the \g{implementation}'s chosen output backend does not match, the instruction is ignored. \\

Users should note that basic Markless parsing rules such as backslash escapes still apply for the \inline$content$, so the content is not copied directly 1:1 from the source text to the output document. \\

\begin{examples}
  \example{! raw latex \\\\textit\{Hello\}}{\textit{Hello}}
  \example{! raw html <blink>Hello</blink>}{}
\end{examples}

%%% Local Variables:
%%% mode: latex
%%% TeX-master: "markless"
%%% TeX-engine: luatex
%%% TeX-command-extra-options: "-shell-escape"
%%% End:

%%% Local Variables:
%%% mode: latex
%%% TeX-master: "markless"
%%% TeX-engine: luatex
%%% TeX-command-extra-options: "-shell-escape"
%%% End:
\input{5.8-comment.tex}
\definesubsection{Embed}
\begin{identifier}{embed}
\[ <type .*> <target ![,]*>(, *<parameter ![,]*>)*( *\])?
\end{identifier}
\definetextualcomponent{embed}{display: block;target: target} \\

The embed is a \g{singular line directive}. The content of the \inline$type$ binding determines the embed's type, and the \inline$parameter$ bindings determine the embed's parameters. The style of the \g{resulting textual component} is dynamically dependant on the given type. \\

Unless the \ident{embed-property-width} or \ident{embed-property-height} parameters are present, the size of the embed \g{resulting textual component} is constrained to be smaller than the width and height of the \g{page} it is output to while preserving the embed content's aspect ratio. If the \g{page} has no width or height, or the embed content's dimensions are smaller than both of those, then the embed content is sized to its own dimensions. The \g{resulting textual component} must not be split across multiple \gpl{page}.\\

\definesubsubsection{Embed Types}
An \g{implementation} must at least support the types specified in this section if permitted by the output \g{document}, but may add additional options the implications of which are completely \g{implementation dependant}. If the output \g{document} does not support a particular type, a \comp{paragraph} containing a single \comp{url} \g{textual component} is outputted with its target set to the \inline$target$ \g{binding}'s value, its content set to either the content of the \inline$description$ option if present, or the \inline$target$ \g{binding}'s value if not, and a \g{warning} is \glink{signalling}{signalled}. If the \g{implementation} does not support the requested type at all, an \g{error} is \glink{signalling}{signalled}. \\

\definesubsubsubsection{Image}
\defineidentifier{embed-type-image}{image}
\definestyle{embed-type-image}{interaction: image} \\

Embeds the image pointed to by the \inline$target$ into the document. The supported image formats are \g{implementation dependant}. If the format of the target is not supported by the \g{implementation}, the \g{directive} is treated as if it were given an unknown type. \\

\begin{examples}
  \example{[ image assets/markless-logo.png ]}{\includegraphics[scale=0.5]{assets/markless-logo}}
\end{examples}

%%% Local Variables:
%%% mode: latex
%%% TeX-master: "markless"
%%% TeX-engine: luatex
%%% TeX-command-extra-options: "-shell-escape"
%%% End:

\input{5.9.1.2-video.tex}
\input{5.9.1.3-audio.tex}
\input{5.9.1.4-source.tex}

%%% Local Variables:
%%% mode: latex
%%% TeX-master: "markless"
%%% TeX-engine: luatex
%%% TeX-command-extra-options: "-shell-escape"
%%% End:

\definesubsubsection{Embed Parameters}
The parameters are processed in the order they are given and can effect both the content of the \g{resulting textual component} as well as its \g{style}. A parameter may also affect the processing of parameters after it. Two general types of parameters are defined: flag parameters and value parameters. Flag parameters are single parameters that add or remove an attribute from the \g{resulting textual component}'s \g{style}. Value parameters add an attribute whose value is determined by the parameter following the current one. The following parameter is then skipped over and thus not processed. \\

An \g{implementation} must at least support the parameters specified in this section if permitted by the output \g{document}, but may add additional parameters the implications of which are completely \g{implementation dependant}. If the output \g{document} does not support a particular parameter, or an unknown parameter is given, a \g{warning} is \glink{signalling}{signalled}. \\

\input{5.9.2.1-float.tex}
\definesubsubsubsection{Width}
\defineidentifier{embed-property-width}{width (<pixels ~n+px>|<percent ~n+\%>)}
\definestyle{embed-property-width}{width: size} \\

Causes the embed content's width to be fixed to the specified size. The size can be given in either \inline$pixels$ or \inline$percent$ where \inline$pixels$ will set the width to be the exact amount of pixels given if the document is viewed at its native resolution. \inline$percent$ will scale the width to the given percentage of the width of the \g{document}. If the \g{document} should not have a width, the \inline$percent$ specification does nothing. Unless the \ident{embed-property-height} is also specified, the embed content's aspect ratio must be preserved. \\

\begin{examples}
  \example{[ image assets/markless-logo.png, width 50px ]}{\includegraphics[width=50px]{assets/markless-logo}}
\end{examples}

%%% Local Variables:
%%% mode: latex
%%% TeX-master: "markless"
%%% TeX-engine: luatex
%%% TeX-command-extra-options: "-shell-escape"
%%% End:

\definesubsubsubsection{Height}
\defineidentifier{embed-property-height}{height (<pixels ~n+px>|<percent ~n+\%>)}
\definestyle{embed-property-height}{height: size} \\

Causes the embed content's height to be fixed to the specified size. The size can be given in either \inline$pixels$ or \inline$percent$ where \inline$pixels$ will set the height to be the exact amount of pixels given if the document is viewed at its native resolution. \inline$percent$ will scale the height to the given percentage of the height of the \g{document}. If the \g{document} should not have a height, the \inline$percent$ specification does nothing. Unless the \ident{embed-property-width} is also specified, the embed content's aspect ratio must be preserved. \\

\begin{examples}
  \example{[ image markless-logo.png, width 50px, height 100px ]}{\includegraphics[width=50px,height=100px]{markless-logo}}
\end{examples}

%%% Local Variables:
%%% mode: latex
%%% TeX-master: "markless"
%%% TeX-engine: luatex
%%% TeX-command-extra-options: "-shell-escape"
%%% End:

\input{5.9.2.4-label.tex}
\definesubsubsubsection{Caption}
\defineidentifier{embed-property-caption}{caption <content .*>} 
\definestyle{embed-property-caption}{interaction: image} \\

Causes a \g{textual component} to be output either alongside or within the embed's \g{textual component}. The text may contain any \gpl{inline directive}. The \g{text} held by the \g{content binding} is outputted to this additional \g{resulting textual component}. The text can only contain \gpl{inline directive}.\\

\begin{examples}
  \longexample{[ image assets/markless-logo.png, caption The //Markless// logo image. ]}{\begin{minipage}{.98\linewidth}
      \begin{center}
        \includegraphics[height=2cm]{assets/markless-logo}
      \end{center}
      \vspace{-1cm}
      \captionof{figure}{The \textit{Markless} logo image.}
    \end{minipage}}
\end{examples}

%%% Local Variables:
%%% mode: latex
%%% TeX-master: "markless"
%%% TeX-engine: luatex
%%% TeX-command-extra-options: "-shell-escape"
%%% End:

\definesubsubsubsection{Description}
\defineidentifier{embed-property-description}{description <content .*>} \\

A textual description of the embed's content for use when the embed content cannot be displayed, or when the \g{user} employs a reader or other form of aid system that relies entirely on textual representation.

%%% Local Variables:
%%% mode: latex
%%% TeX-master: "markless"
%%% TeX-engine: luatex
%%% TeX-command-extra-options: "-shell-escape"
%%% End:

%%% Local Variables:
%%% mode: latex
%%% TeX-master: "markless"
%%% TeX-engine: luatex
%%% TeX-command-extra-options: "-shell-escape"
%%% End:


%%% Local Variables:
%%% mode: latex
%%% TeX-master: "markless"
%%% TeX-engine: luatex
%%% TeX-command-extra-options: "-shell-escape"
%%% End:

\input{5.10-footnote.tex}
\definesubsection{Alignment}
\begin{identifier}{Left Align}
  \|\<<content .*>
\end{identifier}
\begin{identifier}{Right Align}
  \|\><content .*>
\end{identifier}
\begin{identifier}{Center}
  \>\<<content .*>
\end{identifier}
\begin{identifier}{Justify}
  \|\|<content .*>
\end{identifier}
\definetextualcomponent{left align}{text-align: left}
\definetextualcomponent{right align}{text-align: right}
\definetextualcomponent{center}{text-align: center}
\definetextualcomponent{justify}{text-align: justify} \\

All alignment directives are \g{spanning line directive} that change the alignment of a body of \g{text}. The alignment content can contain any \g{directive} with the condition that the \gpl{directive} are matched against the \g{text} of the \g{resulting textual component}. \\

\begin{examples}
  \begin{examplesource}
    |< Left
    >< Center
    |> Right
  \end{examplesource}
  \begin{exampleoutput}
    \begin{minipage}{0.5\textwidth}
      \begin{flushleft}Left\end{flushleft}
      \begin{center}Center\end{center}
      \begin{flushright}Right\end{flushright}
    \end{minipage}
  \end{exampleoutput}
\end{examples}

%%% Local Variables:
%%% mode: latex
%%% TeX-master: "markless"
%%% TeX-engine: luatex
%%% TeX-command-extra-options: "-shell-escape"
%%% End:


%%% Local Variables:
%%% mode: latex
%%% TeX-master: "markless"
%%% TeX-engine: luatex
%%% TeX-command-extra-options: "-shell-escape"
%%% End:
